\documentclass[10pt,a4paper]{amsart}
\usepackage[margin=19mm]{geometry}
\usepackage{amsmath,amssymb,mathtools}
\usepackage[scr=rsfs,cal=euler]{mathalfa}
\usepackage{xfrac}
\usepackage[unicode,hidelinks,bookmarks=false]{hyperref}
\newcommand{\ie}{i.e.\ }
\newcommand{\cf}{cf.\ }
\newcommand{\resp}{respectively\ }
\newcommand{\Subsection}{\subsection}
\providecommand{\nobreakdash}{\nobreak\hbox{-}}
\setlength{\emergencystretch}{2em}
\multlinegap0pt
\usepackage{bm}
\usepackage{bbm}
\usepackage{dsfont}
\usepackage{xspace}
\usepackage{cancel}
\usepackage{mathtools}
\usepackage{amsthm}
\makeatletter
\@ifundefined{theorem}{\newtheorem{theorem}{Theorem}}{}
\makeatother
\renewcommand{\L}{\mathcal{L}}
\renewcommand{\P}{\mathcal{P}}
\title[Transparent Semantic Spaces: A Categorical Approach to Expla]{Transparent Semantic Spaces: A Categorical Approach to Explainable Word Embeddings}
\author{Ares Fabregat-Hern\'andez (1, 2), Javier Palanca (1), Vicent Botti (1, 3) ((1) Valencian Research Institute for Artificial Intelligence (VRAIN) Universitat Polit\`ecnica de Val\`encia (2) Universidad Internacional de Valencia (VIU) (3) valgrAI (Valencian Graduate School, Research Network of Artificial Intelligence))}
\date{Source: arXiv:2508.20701v1 (CC BY 4.0)}
\begin{document}
\maketitle

\noindent\emph{Source and licence.} Transparent Semantic Spaces: A Categorical Approach to Explainable Word Embeddings. Version arXiv:2508.20701v1 (CC BY 4.0), distributed under the Creative Commons Attribution 4.0 International licence, as recorded in the arXiv licence metadata for this exact version and shown on the abstract page https://arxiv.org/abs/2508.20701v1. Full rights notice: licenses/arxiv-2508.20701-CC-BY-4.0.txt.\par\smallskip

\noindent\textit{Editorial scope.} Counted unit: the Theorem whose source label is \texttt{thm:MonoidalCatA} in the article (source0.tex lines 1030--1032, source label \texttt{thm:MonoidalCatA}), delivered together with the article's own complete original proof of it (source0.tex lines 1034--1058, one \texttt{proof} environment of 25 lines). No mathematics is written, repaired or re-derived: statement and proof are the article's own text line for line. Required context retained uncounted: thm:MonoidalCat (lines 525--527). Every definition, display and label of those lines is retained unchanged. Omitted from this package and not redistributed: 1--524, 528--1029, 1033--1033, 1059--1242. Nothing inside a retained excerpt is omitted or altered: every retained body is delivered exactly as the article's lines are written, so no omission receipt of any kind accompanies this package. Citations: the retained excerpts carry no citation command at all, so no bibliography entry of the article is transcribed and no reference is redistributed. Reference adaptation: none was needed and none was made; every reference inside the delivered unit resolves inside it, so no unresolved reference or citation is produced. Printed slips kept literal: no printed word, symbol, hypothesis or number of the counted statement or of its proof was altered. Layout and class: the article's own class is \texttt{elsarticle} with packages including amsthm, amsmath, bbm, amsfonts, amssymb, mathtools, graphicx, natbib, caption, subcaption, dutchcal, lmodern, tikz-cd, hyperref; the wrapper is the lane's pinned \texttt{amsart} editorial wrapper at 10pt on A4, which restates the theorem-like environments the retained text needs and adds the article's own macro definitions that the retained text needs (\texttt{\textbackslash L}, \texttt{\textbackslash P}), with extra wrapper packages (bm, bbm, dsfont, xspace, cancel, mathtools). The delivered numbering is the unit's own. Assets: no retained span contains a figure environment, a graphics-inclusion command, a PDF-page inclusion, a table or a TikZ picture; no image file is copied into this package, the package declares no asset dependency and no image file is redistributed with it. Licence: version arXiv:2508.20701v1 of this article is distributed under the Creative Commons Attribution 4.0 International licence, as recorded in the arXiv licence metadata for this exact version and shown on the abstract page https://arxiv.org/abs/2508.20701v1. A full rights notice is delivered at licenses/arxiv-2508.20701-CC-BY-4.0.txt. Characters: every retained excerpt is pure ASCII, so the delivered engine, XeLaTeX with its default Latin Modern text font, reports no missing character.
\begin{theorem}\label{thm:MonoidalCat}
	The category $(\P_{T},\otimes, \L_{T}^{0})$ is indeed a monoidal category.
\end{theorem}

	\begin{theorem}[Theorem \ref{thm:MonoidalCat}]\label{thm:MonoidalCatA}
		The category $(\P_{T},\otimes, \L_{T}^{0})$ is indeed a monoidal category.
	\end{theorem}

	\begin{proof}
		There are two things to prove. The first is to show that $\P_{T}$ is a category with the operations defined above. The second is that $(\P_{T},\otimes, \L_{T}^{0})$ satisfies all the axioms of a monoidal category.
		
		Let's start with the first part. To show that $\P_{T}$ is a category we have to show that given two morphisms $f\colon X to Y$ and $g\colon Y\to Z$ their composition is also a morphism $h = g\circ f\colon X\to Z$ in $\P_{T}$. Specifically, we have to check that every entry of the product matrix is a positive number between $0$ and $1$. If the two morphisms are row-stochastic matrices, each entry of the matrix $h$ can be expressed as a sum
		
		\begin{equation}\label{eq:def:compmorphA}
			h_{ij}=\sum_{k} f_{ik}g_{kj}\leq \sum_{k} g_{kj}\leq 1
		\end{equation}
		
		since $f$ and $g$ are row-stochastic. Thus, their composition is again a morphism in $\P_{T}$\footnote{In fact, the product of row-stochastic matrices is again row-stochastic.}. 
		
		
		Now, assume one of them is a probabilistic matrix, let's say $f$. Then, each entry of the product is weighted down by the inverse of the ceiling function of the number of the sum of the transition probabilities. This means that equation $\eqref{eq:def:compmorphA}$ becomes
		
		\begin{equation}\label{eq:def:compmorphceilA}
			h_{ij}= \frac{1}{\left\lceil\sum_{j}f_{ij}\right\rceil}\sum_{k} f_{ik}g_{kj} = \sum_{k} \frac{f_{ik}}{\left\lceil\sum_{j}f_{ij}\right\rceil} g_{kj}\leq 1.
		\end{equation}
		Thus, their composition is again a morphism in $\P_{T}$. Equations \eqref{eq:def:compmorphA} and \eqref{eq:def:compmorphceilA} imply the associativity of the composition of morphisms in $\P_{T}$: Finally, the identity morphism for each object $X$ in $\P_{T}$ is the identity matrix. Hence, $\P_{T}$ is a well-defined category.
		
		To show that the tuple $(\P_{T},\otimes, \L_{T}^{0})$ is a monoidal category we need to check that the tensor product we have defined satisfies associativity, that there are isomorphisms between the objects $X\otimes \L_{T}^{0}\cong X$ and $\L_{T}^{0}\otimes X\cong X$ and that the triangle and pentagon identities are satisfied.
		
		First of all, since the tensor product is just a concatenation of expressions of $T$ it is clear that the tensor product is associative. The $0$-th graded piece acts as concatenating $<eos>$ which is as concatenating nothing. Hence, it is the identity and we have $X\otimes \L_{T}^{0}\cong \L_{T}^{0}\otimes X\cong X$ given by $x\otimes <eos> = x = <eos>\otimes x$ for all $x\in X$.
		
		Lastly, the triangle and pentagon identities follow from the definition of the tensor product, the associativity condition, and the identity $x\otimes <eos> = x = <eos>\otimes x$. Thus, the tuple $(\P_{T},\otimes, \L_{T}^{0})$ is a monoidal category.
	\end{proof}

\end{document}
